\documentclass[a4paper,12pt]{article}

\usepackage[T1]{fontenc}
\usepackage[italian]{babel}
\usepackage{graphicx}
\usepackage{geometry}
\usepackage{tabularx}
\usepackage[table]{xcolor}
\usepackage[most]{tcolorbox}
\usepackage{fancyhdr}
\usepackage[hidelinks]{hyperref}
\usepackage{tikz}
\usepackage{pgf-pie}

\newcommand{\groupmail}{alittlebyte.swe@gmail.com}
\newcommand{\grouplogo}{../../.github/site-src/assets/logo_alittlebyte.png}
\newcommand{\groupsymbol}{../../.github/site-src/assets/solo_logo.png}

\geometry{
    left=2.5cm,
    right=2.5cm,
    top=2.5cm,
    bottom=2.5cm,
    headheight=331pt,
    includeheadfoot
}

\pagestyle{fancy}
\fancyhf{}

\renewcommand{\headrulewidth}{0pt}
\fancyhead{}

\renewcommand{\footrulewidth}{0.4pt}
\fancyfoot[L]{\includegraphics[height=0.6cm]{\groupsymbol}\hspace{0.45em}aLittleByte}
\fancyfoot[R]{Analisi dei Capitolati}

\fancypagestyle{firstpage}{
    \fancyhf{}
    \renewcommand{\headrulewidth}{0pt}
    \renewcommand{\footrulewidth}{0.4pt}
    \fancyhead[C]{%
        \makebox[\textwidth][c]{%
            \begin{tabular}{c}
                \includegraphics[height=10cm]{\grouplogo}\\[1ex]
                \small\texttt{\groupmail}
            \end{tabular}%
        }
    }
    \fancyfoot[L]{\includegraphics[height=0.6cm]{\groupsymbol}\hspace{0.45em}aLittleByte}
    \fancyfoot[R]{Dichiarazione Impegni}
}

\begin{document}

\thispagestyle{firstpage}

\vspace*{0.5cm}

\begin{center}
    {\LARGE \textbf{Dichiarazione Impegni}}
\end{center}

\vspace{1.5cm}

\begin{center}
    \renewcommand{\arraystretch}{1.3}
    \begin{tcolorbox}[
        width=0.92\textwidth,
        colback=gray!6,
        colframe=gray!45,
        boxrule=0.5pt,
        arc=2mm,
        boxsep=0pt,
        left=8pt, right=8pt, top=8pt, bottom=8pt
    ]
        \begin{tabularx}{\linewidth}{>{\bfseries}p{3.5cm} X}
            Gruppo      & aLittleByte (gruppo 19) \\
            Redattore   & Angela Maule, Eleonora Bellet \\
            Verificatori & Alex Oloni, Vinicius Yuri Kaneda Fernandes \\
            Destinatari & Prof. Tullio Vardanega, Prof. Riccardo Cardin \\
            Data        & 2026-03-19
        \end{tabularx}
    \end{tcolorbox}
\end{center}

\newpage
\fancyfoot[R]{Dichiarazione Impegni\hspace{0.8em}}
\newgeometry{left=2.5cm,right=2.5cm,top=2.5cm,bottom=2.5cm,includefoot}

\begin{center}
    {\Large \textbf{Registro delle Modifiche}}
\end{center}

\vspace{0.35cm}

\renewcommand{\arraystretch}{1.35}
\rowcolors{2}{gray!8}{white}
\begin{center}
    {\footnotesize
    \begin{tabularx}{\textwidth}{|l|l|X|p{2cm}|p{2cm}|p{2cm}|}
        \hline
        \rowcolor{gray!20}
        \textbf{Versione} & \textbf{Data} & \textbf{Descrizione} & \textbf{Redattore} & \textbf{Verificatore} & \textbf{Approvatore} \\
         \hline 
        1.0.0 & 2026-03-19 & Approvazione documento & - & - & D. Oloni, V. Y. Kaneda Fernandes\\
        \hline 
        0.3.0 & 2026-03-17 & Stesura sezione 3 (Preventivo costi) & A. Maule & D. Oloni & - \\
        \hline 
        0.2.0 & 2026-03-10 & Stesura sezione 2 (Impegno orario e Ruoli) & A. Maule, E. Bellet & V. Y. Kaneda Fernandes & - \\
        \hline
        0.1.0 & 2026-03-09 & Struttura del documento e Introduzione & E. Bellet & D. Oloni & - \\
        \hline 
    \end{tabularx}
    }
\end{center}

\newpage
\pagenumbering{arabic}
\setcounter{page}{1}
\fancyfoot[R]{Dichiarazione Impegni\hspace{0.8em}\thepage}

\begin{center} 
    {\Large \textbf{Indice}}
\end{center}

\vspace{0.6cm}

\tableofcontents

\newpage

\section{Introduzione}
\subsection{Finalità del documento}

Con il seguente documento, il gruppo aLittleByte (Gruppo 19) di Ingegneria del Software dichiara l'impegno orario preventivato per l'implementazione del capitolato C5, proposto da Eggon-Nexum, confermando la propria candidatura per la sua realizzazione.
Nel documento sono riportati: la descrizione dei ruoli previsti per lo svolgimento del progetto, il monte ore assegnato a ciascun componente e il preventivo economico totale.

\section{Impegno orario e Ruoli}
\subsection{Distribuzione ore per ruolo}
\begin{table}[h!]
\centering
\renewcommand{\arraystretch}{1.5}
\begin{tabular}{|l|c|c|c|c|c|c|c|}
\hline
\textbf{Componente} & \textbf{Re} & \textbf{Am} & \textbf{An} & \textbf{Pj} & \textbf{Pr} & \textbf{Ve} & \textbf{Totale} \\
\hline
Eleonora Bellet &8  &7  &8  &20  &23  &24  & \textbf{90} \\
\hline
Diego Belluz &8  &7  &9  &21  &23  &22  & \textbf{90} \\
\hline
Angelica Gastaldello &8  &7  &9  &21  &22  &23  & \textbf{90} \\
\hline
Vinicius Yuri Kaneda Fernandes &8  &8  &8  &20  &23  &23  & \textbf{90} \\
\hline
Angela Maule &8  &8  &8  &21  &23  &22 & \textbf{90} \\
\hline
Alex Oloni &8  &8  &9  &20  &23  &23  & \textbf{90} \\
\hline
Leonardo Soligo &8  &7  &9  &20  &23  &23  & \textbf{90} \\
\hline
\textbf{Totale Ore} & \textbf{56} & \textbf{51} & \textbf{60} & \textbf{143} & \textbf{160} & \textbf{160} & \textbf{630} \\
\hline
\end{tabular}
\caption{Ore di ogni componente per ciascun ruolo}
\label{tab:ore_componenti}
\vspace{0.2cm}
\footnotesize{\textbf{Legenda:} \textbf{Re}: Responsabile, \textbf{Am}: Amministratore, \textbf{An}: Analista, \textbf{Pj}: Progettista, \textbf{Pr}: Programmatore, \textbf{Ve}: Verificatore.}
\end{table}
\subsection{Analisi dei Ruoli}
\subsubsection{Responsabile}
Il responsabile coordina le attività e gestisce le comunicazioni con i docenti e l'azienda proponente (Eggon-Nexum). Si occupa della pianificazione temporale, del monitoraggio dell'avanzamento, assicurandosi che la rotazione dei ruoli avvenga in modo corretto. Ha inoltre il compito di definire e far rispettare le \textit{Norme di Progetto} e le linee guida di sviluppo. È una figura presente durante l'intero ciclo di vita del software, ma il suo impegno è particolarmente intensivo nelle fasi iniziali, per poi diminuire progressivamente man mano che il team acquisisce autonomia operativa. Essendo il ruolo economicamente più oneroso, il suo monte ore è stato ottimizzato per garantire efficienza senza gravare eccessivamente sul preventivo finale.

\subsubsection{Amministratore}
Ha il compito di predisporre e preservare l'ecosistema tecnologico di supporto al team, garantendo la piena operatività delle infrastrutture IT e degli strumenti di collaborazione. Promuove l'efficienza del gruppo integrando processi automatizzati, come la validazione continua e l'esecuzione dei test. Trattandosi di una figura focalizzata in larga parte sul setup iniziale e sul monitoraggio dei sistemi, l'impegno orario a preventivo risulta tra i più contenuti dell'intero progetto.

\subsubsection{Progettista}
Ha il compito di definire l'architettura del sistema a partire dai requisiti concordati, selezionando le tecnologie più idonee per la sua realizzazione. Produce inoltre una documentazione chiara ed esaustiva, fondamentale per guidare i programmatori nella successiva fase di codifica. Poiché eventuali imprecisioni progettuali avrebbero gravi ripercussioni sull'intero sviluppo del prodotto, a questo ruolo è stato assegnato un monte ore particolarmente elevato per garantirne la massima cura.

\subsubsection{Analista}
La figura dell'analista è incaricata di studiare e analizzare i requisiti di sistema. Analizza a fondo il capitolato d'appalto, si confronta direttamente con il proponente e traduce le necessità emerse in chiare specifiche e modelli di casi d'uso. Il suo lavoro si concentra soprattutto all'avvio del progetto, ma resta un punto di riferimento fondamentale per adattare e far evolvere i requisiti durante tutto il processo di lavoro.

\subsubsection{Verificatore}
La figura del verificatore ha il compito di certificare la correttezza e la qualità del prodotto finale. È una presenza costante in ogni fase dello sviluppo, impegnata a validare i rilasci e a garantire che tutto il materiale prodotto rispetti i requisiti stabiliti. Le mansioni includono testing, revisione puntuale del codice e controllo di conformità della documentazione rispetto agli standard. 

\subsubsection{Programmatore}
Il programmatore è la figura operativa per eccellenza. Si occupa della stesura del codice in linea con gli standard di settore. Si fa carico di assemblare i componenti del sistema, documentarne il funzionamento e garantirne l'affidabilità fin dalle basi tramite i test unitari. Il suo intervento è richiesto in maniera continuativa per tutta la durata delle fasi di realizzazione e manutenzione.

\subsection{Motivazione suddivisione oraria }
Il preventivo orario complessivo di 630 ore è stato suddiviso assegnando a ciascun membro un impegno di 90 ore produttive. Per assicurare una crescita professionale completa, il gruppo ha pianificato una rotazione dei ruoli ogni due settimane.
Osservando la distribuzione: il nucleo centrale del lavoro è stato posto sulle figure tecniche. In particolare, si è deciso di assegnare un peso esattamente identico alla stesura del codice (Programmatore, 160 ore) e al controllo qualità (Verificatore, 160 ore). Questa specularità nasce dalla volontà di non relegare il testing a una fase secondaria, ma di farlo procedere di pari passo con lo sviluppo.
\\
Una parte importante del preventivo (143 ore) è stata assegnata al Progettista. L'obiettivo è creare da subito una base solida: investire tempo in una buona architettura e in una documentazione chiara facilita il lavoro dei programmatori e riduce il rischio di dover correggere gravi errori in futuro.
\\
Di conseguenza, le ore per l'Analista (60), il Responsabile (56) e l'Amministratore (51) sono state mantenute volutamente più basse. Non si tratta di ruoli meno importanti, ma di una scelta per ottimizzare i costi: poiché il loro lavoro si concentra soprattutto all'inizio del progetto o in attività di controllo periodico, assegnare loro meno ore ci permette di contenere il preventivo e dedicare più risorse allo sviluppo vero e proprio del software.

\section{Preventivo costi}

\subsection{Costo per ruolo}
\begin{table}[h!]
\centering
\renewcommand{\arraystretch}{1.5}
\begin{tabular}{|l|c|c|c|}
\hline
\textbf{Ruolo} & \textbf{Costo Orario (€/h)} & \textbf{Ore} & \textbf{Costo (€)} \\
\hline
Responsabile & 30€/h &56h  &1.680€  \\
\hline
Amministratore & 20€/h &51h  &1.020€  \\
\hline
Analista & 25€/h &60h  &1.500€  \\
\hline
Progettista & 25€/h &143h  &3.575€  \\
\hline
Programmatore & 15€/h &160h  &2.400€  \\
\hline
Verificatore & 15€/h &160h  &2.400€  \\
\hline
\textbf{Totale} & & \textbf{630} & \textbf{12.575€} \\
\hline
\end{tabular}
\caption{Riassunto dei costi derivanti dalle ore assegnate a ciascun ruolo}
\label{tab:costi_ruoli}
\end{table}
\begin{figure}[h]
    \centering
    \begin{tikzpicture}
        \pie[
            text=legend,
            radius=3,
            color={red!60, blue!60, yellow!60, green!60, purple!60, orange!60}
        ]{
            8.9/Responsabile (56 ore),
            8.1/Amministratore (51 ore),
            9.5/Analista (60 ore),
            22.7/Progettista (143 ore),
            25.4/Programmatore (160 ore),
            25.4/Verificatore (160 ore)
        }
    \end{tikzpicture}
    \caption{Ripartizione in percentuale del monte ore totale}
\end{figure}
\subsection{Costo totale preventivato}
Il costo totale previsto per la realizzazione del progetto è di 12.575€.
\newpage
\section{Mitigazione dei rischi}
Per assicurarci di non superare il budget di 12.575€, il gruppo ha pianificato la gestione dei rischi concentrandosi sull'efficienza operativa. Come già citato in precedenza, la scelta di bilanciare esattamente le ore di programmazione e di verifica (160 ore ciascuna) serve a evitare che eventuali errori scoperti in ritardo portino a un aumento imprevisto dei costi nelle fasi finali del progetto.
Inoltre, la rotazione dei ruoli ogni due settimane è supportata da un rigoroso obbligo di documentazione; grazie a ciò ogni membro del gruppo avrà le competenze necessarie per supportare gli altri durante i cambi di turno, garantendo che l'investimento orario rimanga produttivo e aderente al preventivo stabilito. Tale approccio mira a evitare ritardi nella consegna fissata finale del progetto e la necessità di aggiungere ore extra non preventivate.
\section{Consegna}
La data ultima per la consegna del progetto è fissata improrogabilmente al giorno 10 settembre 2026. 

\end{document}